\section{System Model and Problem Formulation}
\label{sec:model}

\subsection{Scheduling environment and decision events}

We consider an edge--fog--cloud infrastructure composed of a set of
heterogeneous computing nodes
\begin{equation}
    \N=\{n_1,n_2,\ldots,n_M\},
\end{equation}
and a set of microservices
\begin{equation}
    \Sset=\{s_1,s_2,\ldots,s_Q\}.
\end{equation}
A microservice may have one or more running instances. Let
$\mathcal{I}$ denote the set of instances and let
$\sigma(i)\in\Sset$ identify the microservice to which instance $i$
belongs.

Scheduling decisions are made at discrete control timestamps
$t\in\mathcal{T}$ separated by an interval $\Delta$. In the trace-based
evaluation, $\Delta=30$ seconds. At timestamp $t$, the scheduler receives
a batch of decision events
\begin{equation}
    \Eset_t=\{e_{t,1},e_{t,2},\ldots,e_{t,B_t}\},
\end{equation}
where
\begin{equation}
    e_{t,r}=(i_{t,r},s_{t,r},n_{t,r},t)
\end{equation}
represents instance $i_{t,r}$ of microservice $s_{t,r}$ currently
located on source node $n_{t,r}$. Multiple events may occur at the same
timestamp. To make the scheduling process deterministic, events within
each batch are processed according to a fixed event identifier.

Each node $n$ exposes a normalized operational state
\begin{equation}
\label{eq:node-state}
    \mathbf{x}_{n,t}
    =
    \begin{bmatrix}
        c_{n,t} & m_{n,t}
    \end{bmatrix},
    \qquad
    c_{n,t},m_{n,t}\in[0,1],
\end{equation}
where $c_{n,t}$ and $m_{n,t}$ denote the observed CPU and memory
utilization at decision time $t$. Only information available at or
before $t$ may be used when constructing a scheduling decision.

In addition to its operational state, each node has a capability profile
\begin{equation}
\label{eq:node-capability}
    \mathbf{a}_n
    =
    \left(
        j_n,\tau_n,\eta_n,\zeta_n
    \right),
\end{equation}
where $j_n$ denotes the node jurisdiction, $\tau_n$ its trust rank,
$\eta_n\in\{0,1\}$ the availability of the required encryption
capability, and $\zeta_n\in\{0,1\}$ the availability of trusted execution
or hardware-backed isolation.

Each microservice $s$ is associated with a requirement profile
\begin{equation}
\label{eq:service-requirement}
    \mathbf{r}_s
    =
    \left(
        J_s,
        \tau_s^{\min},
        \eta_s^{\min},
        \zeta_s^{\min}
    \right),
\end{equation}
where $J_s$ is the set of permitted jurisdictions,
$\tau_s^{\min}$ is the minimum acceptable trust rank, and
$\eta_s^{\min}$ and $\zeta_s^{\min}$ specify whether encryption and
trusted execution are mandatory.

The policy-admissibility predicate is defined as
\begin{equation}
\label{eq:legal}
\begin{split}
    L(s,n)
    ={}&
    \indicator[j_n\in J_s]\,
    \indicator[\tau_n\geq\tau_s^{\min}]\\
    &{}\times
    \indicator[\eta_n\geq\eta_s^{\min}]\,
    \indicator[\zeta_n\geq\zeta_s^{\min}].
\end{split}
\end{equation}
Accordingly, the policy-admissible node set for service $s$ is
\begin{equation}
\label{eq:legal-set}
    \Lset_s
    =
    \{n\in\N:L(s,n)=1\}.
\end{equation}
Nodes outside $\Lset_s$ are excluded before any operational scoring or
ranking is performed.

\subsection{Near-term operational risk}

Current utilization alone does not indicate whether a node is stable,
recovering, or moving towards overload. We therefore define an
operational-risk target over the next control interval.

Let $\rho_c$ and $\rho_m$ denote the CPU and memory operating
boundaries. In this work,
\begin{equation}
    \rho_c=\rho_m=0.8,
\end{equation}
which reserves operational headroom before full saturation. A node is
labelled unsafe at the next control timestamp when either resource
crosses its boundary:
\begin{equation}
\label{eq:unsafe}
    y_{n,t+\Delta}
    =
    \indicator
    \left[
        c_{n,t+\Delta}>\rho_c
        \;\lor\;
        m_{n,t+\Delta}>\rho_m
    \right].
\end{equation}

Let $\mathbf{h}_{n,t}$ denote the temporal feature representation of
node $n$ at time $t$. It is constructed exclusively from current and
past observations and contains the current resource levels, two lagged
states, first- and second-order changes, short-window means and ranges,
recent variation, and projected-headroom features. The temporal-risk
model estimates
\begin{equation}
\label{eq:risk}
    p_{n,t}
    =
    \Pr
    \left(
        y_{n,t+\Delta}=1
        \mid
        \mathbf{h}_{n,t}
    \right).
\end{equation}

A node passes the temporal-risk screen when
\begin{equation}
\label{eq:risk-acceptance}
    p_{n,t}\leq\theta,
\end{equation}
where $\theta$ is selected using calibration data and remains fixed
during evaluation. The resulting temporally admissible set is
\begin{equation}
\label{eq:temporal-set}
    \mathcal{R}_t
    =
    \{n\in\N:p_{n,t}\leq\theta\}.
\end{equation}

\subsection{Calibrated pre-commit guardrail}

The temporal model estimates the probability of near-term threshold
crossing. A separate pre-commit guardrail verifies whether sufficient
resource headroom remains under a transparent one-step projection.

Let $\mathcal{C}$ denote the calibration-transition set. For CPU and
memory, the one-step persistence residuals are defined as
\begin{align}
\label{eq:cpu-residual}
    r^{c}_{n,t}
    &=
    \left|
        c_{n,t+\Delta}-c_{n,t}
    \right|,\\
\label{eq:memory-residual}
    r^{m}_{n,t}
    &=
    \left|
        m_{n,t+\Delta}-m_{n,t}
    \right|.
\end{align}
The calibration margins are obtained from their empirical
95th percentiles:
\begin{align}
\label{eq:q95-cpu}
    q_c^{0.95}
    &=
    Q_{0.95}
    \left(
        \{r^c_{n,t}:(n,t)\in\mathcal{C}\}
    \right),\\
\label{eq:q95-memory}
    q_m^{0.95}
    &=
    Q_{0.95}
    \left(
        \{r^m_{n,t}:(n,t)\in\mathcal{C}\}
    \right).
\end{align}

The corresponding upper resource projections are
\begin{align}
\label{eq:cpu-upper}
    \widehat{c}^{\,\mathrm{up}}_{n,t+\Delta}
    &=
    c_{n,t}+q_c^{0.95},\\
\label{eq:memory-upper}
    \widehat{m}^{\,\mathrm{up}}_{n,t+\Delta}
    &=
    m_{n,t}+q_m^{0.95}.
\end{align}
We define the normalized guardrail utilization as
\begin{equation}
\label{eq:guardrail}
    g_{n,t}
    =
    \max
    \left\{
        \frac{
            \widehat{c}^{\,\mathrm{up}}_{n,t+\Delta}
        }{\rho_c},
        \frac{
            \widehat{m}^{\,\mathrm{up}}_{n,t+\Delta}
        }{\rho_m}
    \right\}.
\end{equation}
A node passes the pre-commit guardrail when
\begin{equation}
\label{eq:q95-executable}
    g_{n,t}\leq 1.
\end{equation}
The associated guardrail-admissible set is
\begin{equation}
\label{eq:guardrail-set}
    \mathcal{G}_t
    =
    \{n\in\N:g_{n,t}\leq1\}.
\end{equation}

The temporal-risk model and the pre-commit guardrail are intentionally
kept separate. The former uses recent resource dynamics to estimate a
probability of near-term threshold crossing, whereas the latter applies
an explicit calibrated headroom condition immediately before commitment.

\subsection{Simultaneous-event and reservation constraints}

Let
\begin{equation}
\label{eq:active-source-set}
    A_t
    =
    \{n_{t,r}:e_{t,r}\in\Eset_t\}
\end{equation}
denote the set of nodes that host an active source event at timestamp
$t$. When an event requires migration, nodes in $A_t$ are not considered
as destinations. This prevents one event in a simultaneous batch from
being migrated onto the source node of another active event.

Events at timestamp $t$ are processed sequentially in their fixed order.
Let $U_{t,r}(n)$ denote the number of destination reservations assigned
to node $n$ before processing event $e_{t,r}$. A node may accept at most
$C$ new reservations within the same timestamp:
\begin{equation}
\label{eq:reservation-capacity}
    U_{t,r}(n)<C.
\end{equation}
After a destination is committed, its reservation count is updated
before the next event in the batch is processed.

The source node of event $e_{t,r}$ is retained when it satisfies all
three admission conditions:
\begin{equation}
\label{eq:retain-source}
\begin{split}
    R(e_{t,r})
    ={}&
    L(s_{t,r},n_{t,r})\\
    &{}\times
    \indicator[p_{n_{t,r},t}\leq\theta]\,
    \indicator[g_{n_{t,r},t}\leq1].
\end{split}
\end{equation}
If $R(e_{t,r})=1$, no migration is required.

Otherwise, the feasible migration set is
\begin{equation}
\label{eq:feasible}
\begin{split}
    \Fset_{t,r}
    =
    \big\{
        n\in\N:
        &\;L(s_{t,r},n)=1,\\
        &\;p_{n,t}\leq\theta,\\
        &\;g_{n,t}\leq1,\\
        &\;n\notin A_t,\\
        &\;U_{t,r}(n)<C
    \big\}.
\end{split}
\end{equation}

Thus, policy admissibility, temporal-risk acceptance, pre-commit
headroom, active-source exclusion, and reservation capacity are enforced
as hard conditions rather than components of a weighted objective. If
$\Fset_{t,r}$ is empty, the scheduler reports the event as infeasible and
does not commit a constraint-violating destination.

\subsection{Fairness-aware shortlisting}

Let $H_{t,r}(n)$ denote the cumulative number of earlier scheduling
decisions assigned to node $n$ before event $e_{t,r}$. The count includes
all decisions made at earlier timestamps and all preceding decisions
within the current timestamp. It is updated after every committed
selection.

Evaluating and ranking every feasible node for every event is
unnecessary when the infrastructure contains thousands of candidates.
CAGE-S therefore constructs a shortlist
\begin{equation}
\label{eq:shortlist}
    K_{t,r}
    =
    \operatorname{TopK}_{k}
    \left(
        \Fset_{t,r};
        H_{t,r}(n),
        g_{n,t},
        \operatorname{id}(n)
    \right),
\end{equation}
where candidates are ordered lexicographically by historical selection
count, guardrail utilization, and a fixed node identifier. The first
$k$ candidates form the shortlist. This step favours less frequently
selected nodes while retaining deterministic behaviour.

If fewer than $k$ feasible nodes exist, all nodes in $\Fset_{t,r}$ are
retained. If no acceptable destination is found in an initially examined
shortlist, the search is expanded to the remaining nodes in
$\Fset_{t,r}$ before infeasibility is reported.

\subsection{Lexicographic placement objective}

For an event whose source cannot be retained, the destination is selected
from $K_{t,r}$ according to
\begin{equation}
\label{eq:lexicographic}
    n_{t,r}^{*}
    =
    \operatorname*{lexmin}_{n\in K_{t,r}}
    \left(
        H_{t,r}(n),
        p_{n,t},
        g_{n,t},
        \operatorname{id}(n)
    \right).
\end{equation}
The historical count is the primary ranking key, followed by predicted
temporal risk and projected guardrail utilization. The fixed node
identifier resolves exact ties and ensures deterministic replay.

The scheduling decision for event $e_{t,r}$ is therefore
\begin{equation}
\label{eq:complete-decision}
    \pi(e_{t,r})
    =
    \begin{cases}
        n_{t,r},
        & R(e_{t,r})=1,\\[2mm]
        n_{t,r}^{*},
        & R(e_{t,r})=0
          \text{ and }
          \Fset_{t,r}\neq\varnothing,\\[2mm]
        \varnothing,
        & \Fset_{t,r}=\varnothing,
    \end{cases}
\end{equation}
where $\varnothing$ denotes a deferred or infeasible decision.

The formulation is lexicographic rather than additive. Mandatory
compliance and executability conditions first determine which actions
are admissible. Historical usage, temporal risk, and projected headroom
then determine the preferred action within that admissible set. As a
result, a favourable operational score cannot compensate for a failed
policy or pre-commit condition.